\documentclass[11pt, a4paper, oneside]{article}
\usepackage[ngerman]{babel}
\usepackage{worksheet}

\begin{document}
	\author{L. Bung}
	\title{Wiederholung: \hspace{10cm} Trigonometrische Funktionen}
	\subject{Mathematik}
	\class{FHR2}
	\maketitle
	
	\task{Eigenschaften trigonometrischer Funktionen}
	
	Geben Sie für die angegebenen Funktionen jeweils Definitions- und Wertebereich, Periode und Amplitude an.
	
	\begin{multicols}{3}
		\begin{enumerate}[label=\alph*)]
			\item $f(x) = 3\sin(x)+1$
			\item $g(x) = -2\cos(2x)-5$
			\item $h(x) = -\sin(\pi x-\frac{\pi}{2})+4$
		\end{enumerate}
	\end{multicols}
	
	
	\task{Symmetrien trigonometrischer Funktionen}
	
	Nutzen Sie Symmetrie und Periode, um alle Nullstellen der Funktion anzugeben.
	
	\begin{multicols}{2}
		\begin{enumerate}[label=\alph*)]
			\item $f(x) = \sin(x)$
			\item $g(x) = \cos(x)$
		\end{enumerate}
	\end{multicols}
	
	
	\task{Einfache trigonometrische Gleichungen}
	
	Bestimmen Sie die Lösungsmenge der Gleichung für das angegebene Intervall.
	
	\begin{multicols}{3}
		\begin{enumerate}[label=\alph*)]
			\item $f(x) = 2\sin(x) = 0.7$ für $x \in [0; 2\pi]$
			\item $g(x) = 2\cos(x)+1 = 0.7$ für $x \in [2\pi;4\pi]$
			\item $h(x) = 3\sin(x)-1 = 1.2$ für $x \in (0; 2\pi)$
		\end{enumerate}
	\end{multicols}
	
	
	\task{Kompliziertere trigonomische Gleichungen}
	
	Bestimmen Sie die Lösungsmenge der Gleichung für das angegebene Intervall.
	
	\begin{multicols}{3}
		\begin{enumerate}[label=\alph*)]
			\item $f(x) = \sin(3x) = 0$ für $x \in [0; 2\pi]$
			\item $g(x) = 2\cos(2x) = -1$ für $x \in (0; 2\pi]$
			\item $h(x) = 2\sin(\pi x) -1 = 0.5$ für $x \in [-\pi; \pi]$
		\end{enumerate}
	\end{multicols}
\end{document}